% =========================================================================================
% Chapter 3 -- Preliminary Study and State of the Art
% =========================================================================================
\section{Preliminary Study and State of the Art}
\label{sec:sota}

\subsection{Introduction}

This chapter reviews the work this project builds on and measures itself against. It moves
from the clinical setting to the technical one: what point-of-care ultrasound is used for
in emergency medicine, why automated interpretation of ultrasound is harder than of other
imaging modalities, what has been published on the specific organs addressed here, and what
exists in multimodal clinical decision support. It closes by stating the limitations that
recur across this literature and the position this project takes with respect to them.

\subsection{POCUS in Emergency Medicine}
\label{subsec:sota-pocus}

Point-of-care ultrasound differs from radiological ultrasound in who performs it and why.
It is performed by the treating clinician, at the bedside, to answer a specific binary
question arising from the patient in front of them --- is there fluid in the abdomen, is
the left ventricle contracting, is this dyspnoea cardiogenic --- rather than to produce a
comprehensive report.

That framing has produced a set of standardised protocols. The BLUE protocol
\cite{lichtenstein2008blue} reduces the diagnosis of acute respiratory failure to a short
sequence of chest views, each with a defined interpretation, and reports high diagnostic
accuracy for the principal causes of dyspnoea. The \gls{fast} examination applies the same
logic to injury, searching four windows for free fluid. Focused echocardiography in
undifferentiated shock follows a comparable structure.

Two properties of these protocols matter for what follows. First, they are defined in terms
of \emph{findings} rather than diagnoses: the operator records whether B-lines are present,
not whether the patient has heart failure. Second, they are explicitly bounded --- each
protocol states which questions it answers and, by omission, which it does not. Both
properties are adopted directly by the Ultrasound Agent described in
Section~\ref{subsec:ultrasound-agent}.

\paragraph{Operator dependence is measurable, and larger than is often assumed.} The
systematic review by Andersen et al.\ \cite{andersen2019pocus} quantifies it across
applications. False-positive rates varied from 0.7\% to 33.3\% depending on the type of
examination: 0.7--3.2\% for obstetric scans, 0.5--9.9\% for abdominal, and 4.0--33.3\% for
cardiac. A spread of that magnitude within a single application is the empirical form of
the problem stated in Section~\ref{subsec:problem}.

Reported training volumes vary equally widely --- from 2.3 to 31 hours for focused scans,
and from 4 to as much as 320 hours for multi-anatomical examinations. The review's most
useful finding, however, is a negative one: it identified no overall association between
the amount of training and diagnostic accuracy. What determined accuracy was the type of
examination rather than the hours logged, with focused scans outperforming comprehensive
ones, and competence in some applications reached after only a few hours.

That result points the same way as the protocol design discussed above, and it has a direct
architectural consequence here. A bounded question asked of a single organ is answered
reliably; an open-ended examination is not. It is the reason the Ultrasound Agent is built
as a set of narrow per-organ modules, each declaring what it does and does not cover,
rather than as one general-purpose interpreter of ultrasound.

One qualification should be stated with these figures. The review covers general practice
rather than emergency medicine, and the two settings differ in case mix, in operator
training and in the urgency under which examinations are performed. The mechanism ---
accuracy governed by the boundedness of the question asked --- transfers; the specific
percentages should not be read as emergency-department values.

\subsection{AI for Medical Ultrasound}
\label{subsec:sota-ai-us}

Ultrasound is a harder target for computer vision than computed tomography or magnetic
resonance imaging, for reasons intrinsic to the modality rather than incidental to it.

\paragraph{The image depends on the operator.} A CT volume of a given patient is
essentially reproducible; an ultrasound image is a function of probe position, angle,
pressure, depth setting, gain and patient habitus. Two competent operators produce visibly
different images of the same anatomy, so a model must be robust to a source of variation
that barely exists in other modalities.

\paragraph{Speckle is signal and noise at once.} The granular texture of ultrasound arises
from interference of the returning wave. It is not additive noise that can be filtered away
without loss, and conventional denoising removes diagnostic information along with it.

\paragraph{Several findings are artefacts, not anatomy.} This is the property most often
overlooked. \gls{blines} and \gls{alines} do not correspond to structures in the lung; they
are reverberation artefacts whose presence and distribution carry the diagnostic meaning
\cite{lichtenstein2008blue}. A model trained to segment anatomy is solving the wrong problem
for lung ultrasound, which is why the lung module in this project is a classifier and the
cardiac module is a segmentation network.

\paragraph{There is no canonical orientation.} Manufacturers and operators differ in how the
image is presented, and public datasets inherit those differences without documenting them.
This bears directly on the external-validation gap reported in
Section~\ref{subsec:cardiac-results}.

\subsection{Ultrasound Segmentation and Classification}
\label{subsec:sota-seg-cls}

Two task families cover the work in this project.

\emph{Segmentation} assigns a label to every pixel, and is the appropriate formulation when
the clinical quantity is geometric --- a chamber area, a volume, an ejection fraction. The
encoder--decoder architecture introduced by U-Net \cite{ronneberger2015unet}, in which
skip connections carry high-resolution detail across the bottleneck, remains the basis of
most medical segmentation work. Modern variants typically substitute a stronger pretrained
encoder, such as EfficientNet \cite{tan2019efficientnet}, while keeping the overall shape.

\emph{Classification} assigns labels to a whole image or clip, and is appropriate when the
clinical question is about presence rather than extent. It divides further into
single-label problems, where classes are mutually exclusive, and multi-label problems,
where they are not. The distinction is not cosmetic: a lung clip may show B-lines and
consolidation simultaneously, so the lung module must emit independent probabilities per
finding rather than a distribution over one categorical variable.

\paragraph{A methodological caveat that recurs throughout this literature.} Many published
results report metrics computed over frames or images. Where several frames originate from
one patient or one acquisition, a partition drawn at frame level places near-duplicates on
both sides of it, and the resulting score measures memorisation as much as generalisation.
Such figures are not directly comparable with case-level results. This is treated at length
in Chapter~\ref{sec:data}, because it governs the interpretation of every number reported
in this report.

\subsection{Existing POCUS AI Approaches}
\label{subsec:sota-existing}

\subsubsection{nnU-Net and the CAMUS benchmark}

The \gls{camus} dataset \cite{leclerc2019camus} is the reference resource for cardiac
ultrasound segmentation. It comprises 1{,}000 two-dimensional sequences --- two-chamber and
four-chamber views from 500 patients --- with expert contours of the left ventricular
endocardium, the epicardium and the left atrium at the end-diastolic and end-systolic
instants. Crucially for this project, its patient structure is explicit, which is what
makes correct grouping possible without reconstruction.

The accompanying study established that encoder--decoder architectures of the U-Net family
outperform earlier methods and reach accuracy within expert variability. The dataset is
maintained as an open challenge, and the current leaderboard provides the benchmark against
which the cardiac module of this project is measured. Tab.~\ref{tab:camus} reproduces the
leading entry alongside the intra-observer variability recorded for the reference contours.

\begin{table}[H]
  \centering
  \small
  \begin{tabular}{lcccc}
    \toprule
    & \multicolumn{2}{c}{\textbf{Best reported}} &
      \multicolumn{2}{c}{\textbf{Intra-observer}} \\
    \cmidrule(lr){2-3}\cmidrule(lr){4-5}
    \textbf{Structure} & ED & ES & ED & ES \\
    \midrule
    LV endocardium & 0.952 & 0.935 & 0.945 & 0.930 \\
    LV epicardium  & 0.963 & 0.959 & 0.957 & 0.951 \\
    Left atrium    & 0.902 & 0.935 & --- & --- \\
    \bottomrule
  \end{tabular}
  \caption{Dice coefficients on the CAMUS challenge leaderboard, for the leading
    nnU-Net-based entry, with the intra-observer variability of the reference contours for
    comparison. No intra-observer figure is published for the left atrium.}
  \label{tab:camus}
\end{table}

Two observations matter for later chapters. The leading automated results now equal or
exceed intra-observer agreement on both ventricular structures, so the remaining headroom
on this dataset is small. And performance is consistently lower on the left atrium and on
low-quality or artefact-heavy images --- a pattern the cardiac module of this project
reproduces, and which is discussed in Section~\ref{subsec:cardiac-results}.

nnU-Net \cite{isensee2021nnunet} approaches the problem differently, configuring
preprocessing, architecture and training schedule automatically from properties of the
dataset. Its relevance here is as a benchmark rather than a component: it defines the
performance a well-configured standard pipeline achieves on this data, and the cardiac
results in Section~\ref{subsec:cardiac-results} are reported against it.

\subsubsection{USCL: contrastive pretraining on ultrasound}

Most ultrasound models initialise from ImageNet weights, despite natural photographs
sharing little with speckle texture. USCL \cite{chen2021uscl} addresses this directly,
pretraining on ultrasound video through semi-supervised contrastive learning. The authors
assembled US-4, comprising over 23{,}000 images from four ultrasound video sub-datasets,
and report that a USCL-pretrained backbone reaches over 94\% fine-tuning accuracy on the
POCUS lung dataset against 84\% for ImageNet initialisation.

The motivation is sound and the reported gain is substantial, which is why this project
evaluated it. The measurement obtained here disagreed: on the gallbladder task, USCL
initialisation performed markedly worse than ImageNet initialisation. That result is
reported in Section~\ref{subsec:results-rejected} with its numbers. It is not presented as
a refutation of the original work --- the published claim concerns lung classification on
the dataset the pretraining corpus was drawn from, whereas the task here was multi-class
gallbladder classification on a small number of cases --- but it is the reason the
technique was not adopted, and the discrepancy is stated rather than omitted.

\subsubsection{Auto-US: an ultrasound video diagnosis agent}

Auto-US \cite{yang2025autous} is the closest published work to the present project and
deserves detailed comparison. It is described by its authors as the first agent combining
deep learning with \glspl{llm} for ultrasound video-assisted diagnosis, and is organised
into three stages --- video perception, clinical reasoning and diagnostic decision-making
--- which parallel the perception, representation and reasoning layers of
Fig.~\ref{fig:architecture}.

Its perception component, CTU-Net, is a CNN--Transformer network exploiting temporal,
spatial and frequency information in ultrasound video. It is evaluated on the CUV dataset,
assembled by the authors from open-access sources: 495 videos spanning five disease
categories across three organs. The reported classification accuracy is 86.73\%, and
clinician ratings of the generated diagnostic text exceed 3 out of 5.

The convergence is notable: two independent efforts arrived at the same decomposition of
perception followed by language-model reasoning. The differences are equally instructive
and define where this project contributes.

Auto-US pairs ultrasound video with clinical \emph{text} --- chief complaint, physical
examination and additional information --- rather than with structured measurements. Vital
signs and laboratory values are not among its inputs. The paper reports no confidence
calibration or uncertainty quantification for its predictions, no explicit representation
of missing or unmeasured data, and no mechanism for resolving disagreement between the
video classifier and the language model's reasoning. Those four mechanisms are precisely
what this project treats as central, which places its contribution alongside Auto-US rather
than in competition with it.

The authors themselves make the case for going further, observing that ``relying solely on
ultrasound imaging is insufficient for obtaining accurate diagnostic opinions''
\cite{yang2025autous}. Extending that observation from clinical text to vital signs and
laboratory results, and adding an explicit account of what the models cannot exclude, is
the step taken here.

\subsubsection{The POCUS lung ultrasound dataset}

The lung module uses the public dataset released by Born et al.
\cite{born2021pocus}. It comprises 202 videos and 59 images acquired with convex and linear
probes from 216 patients, covering COVID-19, bacterial pneumonia, non-COVID viral
pneumonia and healthy controls, aggregated from heterogeneous sources.

Its strengths and weaknesses both follow from that construction. The heterogeneity is
representative of real point-of-care acquisition, which is desirable. But provenance is
mixed, acquisition is uncontrolled, and --- decisively for this project --- the release
carries no patient identifier linking the multiple clips belonging to one patient. Case
identity therefore has to be reconstructed before any honest partition can be drawn, which
is the subject of Section~\ref{subsec:data-grouping}. \rmq{The dataset has a published
erratum; check which version the downloaded copy corresponds to and cite accordingly.}

\subsection{Multimodal Clinical Decision Support}
\label{subsec:sota-multimodal}

Combining imaging with the rest of the clinical record is the problem this project is
ultimately about, and it has been studied independently of ultrasound. The systematic
review by Huang et al.\ \cite{huang2020fusion} surveys deep-learning systems that fuse
medical imaging with electronic health record data and provides the vocabulary used here.

It distinguishes three strategies. \emph{Early fusion} concatenates the input modalities
into a single feature vector before any model sees them. \emph{Joint fusion} combines
learned intermediate representations, with the loss propagated back into each
modality-specific encoder during training. \emph{Late fusion} trains a separate model per
modality and aggregates their predictions afterwards.

Two of the review's findings matter for the design adopted here. The first is that fusion
works: across the studies surveyed, multimodal models improved accuracy by 1.2--27.7\% and
AUROC by 0.02--0.16 over single-modality baselines, though no single fusion strategy was
consistently optimal across domains. The second is a distributional observation --- of the
17 studies reviewed, 11 (65\%) used early fusion and only 3 (18\%) used late fusion. The
field's default is to fuse early, inside one model.

\paragraph{The property that decides the question for an emergency system.} The same review
notes that late fusion is uniquely able to produce a prediction when not all modalities are
present, because each modality has its own model and the aggregation step can operate over
whichever predictions exist. Early and joint fusion architectures expect a complete input
vector, and have no defined behaviour when part of it is absent.

For a system evaluating an emergency patient this is not a secondary consideration. Missing
data is the normal condition, not the exception: the troponin has not come back, the
history is unobtainable, no ultrasound of that organ was performed. An architecture whose
behaviour is undefined when a modality is missing is undefined most of the time.

The evidence that late fusion also performs competitively is direct. In a case study on
pulmonary embolism detection from CT together with EHR data, Huang et al.\
\cite{huang2020pe} report a late-fusion model reaching an AUROC of 0.947, exceeding the
EHR-only model by 0.036 and the imaging-only model by 0.156. Late fusion was the
best-performing strategy tested, not a concession made for the sake of robustness.

\paragraph{Where this project goes further, and on what basis.} The architecture described
in Section~\ref{subsec:architecture} is late fusion in the sense above --- separate models
per modality, combined afterwards --- but the combination is not an aggregation function
over probabilities. The perception outputs are written into a structured clinical state
that records confidence, evidence grade, model scope and what was not measured, and a
reasoning layer consumes that state.

The motivation is attribution. Under early or joint fusion a single confidence emerges from
the combined model, and a clinician reading it cannot tell whether the imaging or the vital
signs produced it; under the design used here, each finding retains the confidence of the
module that produced it, and disagreement between modules remains visible instead of being
averaged away.

This last point is offered as a design argument rather than as an established result. The
systematic review cited above does not compare fusion strategies with respect to
interpretability, and no measurement in this project isolates the benefit of attribution
from the other properties of the architecture. It is a reason for the design, not evidence
for it.

\subsection{Large Language Models and Retrieval-Augmented Generation}
\label{subsec:sota-llm}

Large language models can integrate heterogeneous evidence and produce structured reasoning
over it, which is the capability the reasoning layer of this system requires. Domain-adapted
medical reasoning models such as HuatuoGPT-o1 \cite{chen2024huatuogpto1} extend this to
clinical material specifically.

Retrieval-augmented generation \cite{lewis2020rag} supplies a model with passages retrieved
from a trusted corpus at inference time, so that its output is grounded in citable sources
rather than in its parameters alone. For clinical use the attraction is less the additional
knowledge than the traceability: an assertion accompanied by the passage it came from can be
checked.

\paragraph{The failure mode that motivates the safety layer.} A fluent model asked to reason
about a patient will produce a fluent answer whether or not the evidence supports one. Left
unconstrained it will cite a laboratory value that was never measured, or treat the absence
of a finding as evidence of health. Retrieval reduces but does not eliminate this: it
grounds the general knowledge while leaving the patient-specific assertions ungrounded.
Neither retrieval nor domain adaptation constitutes a guarantee, which is why in this
project the safety-critical determinations are computed deterministically and the model's
output is validated against the structured state afterwards
(Chapter~\ref{sec:reasoning}).

\subsection{Limitations of Existing Approaches}
\label{subsec:sota-limits}

Five limitations recur across the work reviewed above. They are stated as properties of the
literature, not as criticism of individual papers, most of which address a narrower question
than the one posed here.

\begin{enumerate}
  \setlength\itemsep{0.25em}
  \item \textbf{One organ, one question.} Models are developed and evaluated on a single
    organ. An undifferentiated emergency patient does not present a single question, and
    nothing in these systems combines their outputs.
  \item \textbf{Evaluation below the level of the patient.} Metrics are frequently computed
    over frames or images, inflating measured performance wherever several derive from one
    patient.
  \item \textbf{Uncalibrated confidence reported as probability.} A softmax output is
    routinely presented as a probability without being checked against observed frequency
    \cite{guo2017calibration}. A clinician reading 0.8 as ``correct four times in five''
    is entitled to be wrong.
  \item \textbf{No representation of absent evidence.} A test never performed and a test
    returning normal are distinct states, and systems that lack the distinction implicitly
    treat the first as the second.
  \item \textbf{No defined behaviour when sources disagree.} Where multiple signals are
    available, disagreement is typically resolved by whichever model is more confident,
    rather than being surfaced as the clinically important event it is.
\end{enumerate}

\subsection{Positioning of Our Approach}
\label{subsec:sota-positioning}

\begin{keystatement}
The individual models in this project are standard architectures and are not claimed as
contributions. The contribution is the layer above them.
\end{keystatement}

The cardiac module is a U-Net with an EfficientNet encoder; the lung and gallbladder
modules are EfficientNet classifiers; the triage model is gradient-boosted trees. Each was
selected because it is a defensible standard choice for its task, not because it is novel,
and Chapter~\ref{sec:results} reports them against published benchmarks rather than
claiming to advance them.

What this project adds is the structure in which they cooperate, addressing the five
limitations above in order: a shared output contract across four modules covering both
imaging and structured data; case-level evaluation throughout; a calibrated confidence per
module, carried through to the reasoning step rather than discarded at the module boundary;
an explicit representation of what was not measured and what each model cannot exclude; and
a deterministic policy that escalates on disagreement instead of resolving it.

Tab.~\ref{tab:positioning} sets this against the work reviewed in this chapter. It compares
\emph{capabilities rather than performance}, and the distinction is deliberate. The systems
reviewed here address different organs, different tasks and different datasets: a Dice
coefficient for cardiac segmentation, a classification accuracy on ultrasound video and a
balanced accuracy on gallbladder cases are not commensurable quantities, and placing them in
a single column would invite a comparison that none of them supports. Measured performance
is reported in Chapter~\ref{sec:results}, where each module is compared against its own
baseline and against the benchmark for its own task.

The table should therefore be read across the rows: \emph{yes} indicates that a system
provides the mechanism, \emph{no} that it is a system which does not, and \emph{n/a} that
the contribution is a dataset, benchmark or training method to which the question does not
apply. The bottom row is the only one carrying \emph{yes} in every column, and that is the
claim this project makes --- not that its models are more accurate than the published ones,
but that they are assembled into a system that knows what it does not know.

\begin{table}[H]
  \centering
  \small
  \setlength{\tabcolsep}{4pt}
  \begin{tabular}{llccccc}
    \toprule
    \textbf{Work} & \textbf{Type} & \textbf{Organs} & \textbf{Non-imaging} &
    \textbf{Calibrated} & \textbf{Absent} & \textbf{Conflict} \\
    & & & \textbf{inputs} & \textbf{confidence} & \textbf{evidence} & \textbf{policy} \\
    \midrule
    CAMUS / nnU-Net \cite{leclerc2019camus, isensee2021nnunet}
      & benchmark   & 1 & none & n/a & n/a & n/a \\
    USCL \cite{chen2021uscl}
      & pretraining & 1 & none & n/a & n/a & n/a \\
    POCUS lung \cite{born2021pocus}
      & dataset     & 1 & none & n/a & n/a & n/a \\
    Auto-US \cite{yang2025autous}
      & agent       & 3 & clinical text & no & no & no \\
    \addlinespace[3pt]
    \textbf{This work}
      & agent       & 3 + triage & vitals, labs & yes & yes & yes \\
    \bottomrule
  \end{tabular}
  \caption{Positioning against the reviewed work. ``n/a'' marks a mechanism that does not
    apply to the kind of contribution: a dataset or a pretraining method is not expected to
    provide a conflict policy. Only Auto-US and the present work are complete systems, and
    only those two rows are directly comparable.}
  \label{tab:positioning}
\end{table}

\subsection{Conclusion}

The literature provides strong single-organ models, a reference benchmark for cardiac
segmentation, public data for lung ultrasound, and --- in Auto-US --- an independent
confirmation that perception followed by language-model reasoning is a reasonable
decomposition. What it does not provide is a system that combines imaging with the rest of
the emergency record while keeping an honest account of its own reliability.

The next chapter turns to the data, and to the question that had to be answered before any
model could be trained: what a single independent observation actually is in each of these
datasets.
